\subsection{Hamiltonian and energy origin}
Reduced pairing models provide a useful algebraic test case in the
Richardson--Gaudin literature \citep{Richardson1963,Dukelsky2004}.
The following two-level calculation is derived directly and needs no general
integrability theorem. It is a toy Hamiltonian, not an ab initio model of
$\mathrm H_2$. No molecular integrals or nuclear repulsion are specified.

Order the four fermionic modes as
$(1\uparrow,1\downarrow,2\uparrow,2\downarrow)$ and restrict to two
electrons. Put
\[
 P_p^\dagger=a_{p\uparrow}^\dagger a_{p\downarrow}^\dagger,
 \quad P_p=a_{p\downarrow}a_{p\uparrow},\quad
 n_{p\sigma}=a_{p\sigma}^\dagger a_{p\sigma},\quad p=1,2.
\]
For real $\varepsilon_1,\varepsilon_2$ and attractive coupling $g\ge0$,
define
\begin{align}
 H_z&=\sum_{p=1}^2\varepsilon_p
       (n_{p\uparrow}+n_{p\downarrow}-1)
       -g\sum_{p,q=1}^2P_p^\dagger P_q,
       \label{eq:pair-z}\\
 H_N&=\sum_{p=1}^2\varepsilon_p
       (n_{p\uparrow}+n_{p\downarrow})
       -g\sum_{p,q=1}^2P_p^\dagger P_q
     =H_z+(\varepsilon_1+\varepsilon_2)I.
       \label{eq:pair-n}
\end{align}
The two conventions differ only by an energy origin, but they must not be
mixed in a secular equation or a numerical comparison.

\begin{proposition}[Complete spectrum in the two-electron sector]
\label{prop:pairing}
Let $d=\varepsilon_1-\varepsilon_2$ and
$\rho=\sqrt{d^2+g^2}$. In the paired basis
$|1\overline1\rangle=P_1^\dagger|0\rangle$,
$|2\overline2\rangle=P_2^\dagger|0\rangle$, the matrix of $H_z$ is
\begin{equation}
 H_z^{\mathrm{pair}}=
 \begin{pmatrix}
 \varepsilon_1-\varepsilon_2-g&-g\\
 -g&-\varepsilon_1+\varepsilon_2-g
 \end{pmatrix}.
 \label{eq:pair-block}
\end{equation}
The complete six-dimensional sector has spectrum, with multiplicity,
\begin{equation}
 \{-g-\rho,\ -g+\rho,\ 0,\ 0,\ 0,\ 0\}.
 \label{eq:pair-spectrum}
\end{equation}
In particular its lowest energy is $E_z^-=-g-\rho$.
The paired energies of $H_N$ are
\begin{equation}
 E_N^\pm=\varepsilon_1+\varepsilon_2-g\pm\rho,
 \label{eq:number-roots}
\end{equation}
and its four unpaired energies are $\varepsilon_1+\varepsilon_2$.
\end{proposition}
\begin{proof}
There are $\binom42=6$ two-electron determinants. Two occupy both spins
of a single orbital; the remaining four place one electron in each orbital.
On $|p\overline p\rangle$, $P_q$ is zero unless $q=p$, in which case
it gives the vacuum with positive sign. This follows directly from the CAR:
$a_{p\downarrow}a_{p\uparrow}a_{p\uparrow}^\dagger
 a_{p\downarrow}^\dagger|0\rangle=|0\rangle$.
Subsequent application of $P_s^\dagger$ gives
$|s\overline s\rangle$. Thus $\sum_{s,q}P_s^\dagger P_q$ has every
entry equal to one in the paired basis. In particular the terms $s=q$
contribute $-g$ to each paired diagonal; they cannot be omitted.
The one-body term in $H_z$ is $d$ on $|1\overline1\rangle$ and
$-d$ on $|2\overline2\rangle$. This proves \eqref{eq:pair-block}.

Every pair annihilator kills an unpaired determinant. Its occupations also
give $n_{p\uparrow}+n_{p\downarrow}-1=0$ for each $p$, so its energy
is zero and it is decoupled from the paired space. Finally,
\[
 \det(H_z^{\mathrm{pair}}-\lambda I)
 =(d-g-\lambda)(-d-g-\lambda)-g^2
 =(\lambda+g)^2-d^2-g^2.
\]
Its roots are $-g\pm\rho$. Since $\rho\ge g\ge0$, the lower root
is at most zero and the upper root is at least zero. Comparing with the four
unpaired zeros establishes the ground energy. Adding the scalar in
\eqref{eq:pair-n} shifts every eigenvalue as stated.
\end{proof}

The upper paired root is not generally the first excited level of the full
sector: the unpaired zeros must also be included. When $g>0$, a normalized
ground vector in the paired space is
\[
 c^- =\frac{(g,d+\rho)^T}{\sqrt{g^2+(d+\rho)^2}}.
\]
Indeed the first row of
$(H_z^{\mathrm{pair}}-E_z^-I)(g,d+\rho)^T$ vanishes by cancellation,
and the second is $-g^2+(\rho-d)(\rho+d)=0$.
The denominator is positive for $g>0$. Its orthogonal vector
$c^+=(-c^-_2,c^-_1)^T$ is the upper eigenvector: symmetry makes the
orthogonal one-dimensional subspace invariant, and the trace fixes its
eigenvalue to $-g+\rho$. For $g=0$, determinant basis vectors diagonalize
the block. If also $d=0$, both paired energies and all unpaired energies
coincide at zero. These cases should be handled directly rather than by
dividing by $g$.

\subsection{The one-pair secular equation and its excluded points}
In the number convention,
$H_N^{\mathrm{pair}}=\operatorname{diag}(2\varepsilon_1,
2\varepsilon_2)-g\boldsymbol 1\boldsymbol 1^T$.
An eigenvector $(c_1,c_2)^T$ with energy $E$ obeys
\[
 (2\varepsilon_p-E)c_p=g(c_1+c_2),\qquad p=1,2.
\]
If neither denominator vanishes, a nonzero vector must have
$c_1+c_2\ne0$; otherwise both coefficients would be zero. Dividing and
summing therefore gives
\begin{equation}
 1=g\left(\frac{1}{2\varepsilon_1-E}
          +\frac{1}{2\varepsilon_2-E}\right).
 \label{eq:richardson}
\end{equation}
Conversely, a finite $E$ satisfying this equation away from its poles yields
an eigenvector with coefficients proportional to
$(2\varepsilon_p-E)^{-1}$, by substitution in the two linear equations.
The polynomial form is
\begin{equation}
 (2\varepsilon_1-E)(2\varepsilon_2-E)
 -g\bigl[(2\varepsilon_1-E)+(2\varepsilon_2-E)\bigr]=0.
 \label{eq:pair-polynomial}
\end{equation}
This is the determinant equation and remains meaningful where the rational
form is undefined. It produces exactly \eqref{eq:number-roots}.
At $g=0$ the eigenvalues lie at poles of the rational expression. At
$\varepsilon_1=\varepsilon_2=\varepsilon$ and $g>0$, the upper state
$(1,-1)^T/\sqrt2$ has $E=2\varepsilon$ and zero coefficient sum;
it too is excluded from the rational derivation. A pole is not a zero
residual. Near a pole, the polynomial or matrix eigenvector residual is the
appropriate numerical check, with tolerances scaled to the matrix norm.

\subsection{Relation to the energy certificate}
The paired block is itself a scalar active/complement split with
$a=d-g$, $c=-d-g$, and $b=g$. Its smaller eigenvalue equals the bound
in Theorem~\ref{thm:block}, providing a sharp algebraic check. If the lower
diagonal entry is chosen as the active block and $d\ne0$, the complementary
diagonal lies higher and Corollary~\ref{cor:gap} also applies to this paired
problem. To claim the ground of the full sector one must still account for
the unpaired block; Proposition~\ref{prop:pairing} does so explicitly.

Agreement between a CAR matrix construction, the roots above, and an
independent Jordan--Wigner construction tests this toy algebra. It supplies no
error estimate between a molecular Coulomb Hamiltonian and a reduced pairing
Hamiltonian. Establishing such a reduction would require specified molecular
integrals and a certified bound on the omitted interactions. General claims
about Yangian symmetry, Howe duality, or exact molecular solvability are not
needed for this example and are not established here.
